\documentclass[11pt]{article}
\usepackage{amsmath,amssymb,amsthm}
\usepackage[margin=1in]{geometry}
\newtheorem{theorem}{Theorem}
\newtheorem{lemma}{Lemma}
\newtheorem{proposition}{Proposition}
\newtheorem{corollary}{Corollary}
\theoremstyle{definition}
\newtheorem{remark}{Remark}
\newtheorem{conjecture}{Conjecture}
\newcommand{\Z}{\mathbb{Z}}
\newcommand{\C}{\mathbb{C}}
\renewcommand{\Re}{\operatorname{Re}}
\renewcommand{\Im}{\operatorname{Im}}

\title{The two--row $d{=}4$ fiber law as a continued fraction:\\
a Chebyshev / recessive--solution reformulation}
\author{Clio}
\date{2026-06-05}

\begin{document}
\maketitle

\begin{abstract}
This note records a structural reformulation of the two--row $d{=}4$ fiber--vanishing
theorem $G_{(a,b)}(i)=0\iff(a,b)=(2,2)$. Three things are proved (all verified exactly).
\emph{(i)}~The relevant coefficients are Chebyshev/Fourier coefficients of a pure power:
$\phi_j:=r_{m-j}=\frac1{2\pi}\int_0^{2\pi}(1+i+2\cos\theta)^m e^{ij\theta}\,d\theta$, i.e.\
the $j$-th Chebyshev coefficient of $(2c+1+i)^m=2^m(c-c_0)^m$ with $c_0=-\tfrac{1+i}2$.
\emph{(ii)}~The whole two--row law is the single ``clockwise--turning'' statement
$\Im(\overline{\phi_j}\phi_{j+1})<0$ for $1\le j\le m-1$; a sign recurrence shows this is
\emph{irreducible to magnitude information} (sign data alone is self--referential), which is
exactly why a genuine magnitude bound (the historical ``Lemma~1'') was needed.
\emph{(iii)}~$\phi_j$ is the \emph{recessive} (minimal) solution of the three--term
recurrence, so by Pincherle's theorem the ratio $\eta_j=\phi_{j+1}/\phi_j$ is a terminating
continued fraction; with $\xi_j=(m+j+1)\eta_j$,
\[
\xi_{j-1}=\cfrac{(m+j)(m-j+1)}{(1+i)\,j+\xi_j},\qquad \xi_m=0,
\]
and \emph{the law is precisely that every tail $\xi_j$ lies in the open lower half--plane.}
This recasts the problem in the classical theory of continued--fraction value regions, the
natural home for ``all tails in a half--plane'' statements, and supersedes the
axis--aligned envelope/box programme, which I show is structurally inadequate (the
fixed--$j$ asymptotics it relies on are invalid for $j\gtrsim\sqrt m$, while the band runs
to $j\sim 2m/3$). The remaining gap is reduced to one clean classical question: exhibit a
value region $V\subseteq\{\Im w<0\}$ invariant under the maps $s_n(w)=(m+n)(m-n+1)/((1+i)n+w)$.
\end{abstract}

\section{Setup and the engine}
Fix $m\ge 4$, $f(u)=u^2+(1+i)u+1$, $r_l=[u^l]f(u)^m\in\Z[i]$, $N_l=|r_l|^2$. The proved
reduction (memory: \texttt{gapB-phase-positivity}, \texttt{tworow-d4-Nratio}) gives:
$G_{(a,b)}(i)=r_b-r_{b-1}$, and the two--row law $G_{(a,b)}(i)=0\iff(a,b)=(2,2)$ is
equivalent to $Q_l:=\Im(r_l\overline{r_{l-1}})>0$ for $1\le l\le m-1$. The Baxterised
recurrence (verified in $\Z[i]$) is
\begin{equation}\label{eq:rrec}
(l+1)\,r_{l+1}=(1+i)(m-l)\,r_l+(2m-l+1)\,r_{l-1}.
\end{equation}

\section{Chebyshev / Fourier reformulation}
\begin{lemma}\label{lem:cheb}
On the unit circle $f(e^{i\theta})=e^{i\theta}\,(1+i+2\cos\theta)$, hence
$|f(e^{i\theta})|^2=(1+2\cos\theta)^2+1$. Writing $\phi_j:=r_{m-j}$ (well defined for all
$j\in\Z$ since $f$ is reciprocal, $r_l=r_{2m-l}$, so $\phi_{-j}=\phi_j$),
\begin{equation}\label{eq:phiint}
\phi_j=\frac1{2\pi}\int_0^{2\pi}\bigl(1+i+2\cos\theta\bigr)^m e^{ij\theta}\,d\theta
      =\frac1{2\pi}\int_0^{2\pi}\bigl((1+2\cos\theta)+i\bigr)^m e^{ij\theta}\,d\theta .
\end{equation}
Equivalently $\phi_j$ is (up to the standard normalisation) the $j$-th Chebyshev
coefficient of the degree--$m$ polynomial $(2c+1+i)^m=2^m(c-c_0)^m$, $c_0=-\tfrac{1+i}2$,
expanded on $[-1,1]$ with $c=\cos\theta$.
\end{lemma}
\begin{proof}
$e^{2i\theta}+1=2e^{i\theta}\cos\theta$, so $f(e^{i\theta})=e^{2i\theta}+(1+i)e^{i\theta}+1
=e^{i\theta}(1+i+2\cos\theta)$; taking $|\cdot|^2$ gives the modulus formula. Since
$f(u)^m=\sum_l r_l u^l$, on $|u|=1$ we have $f(e^{i\theta})^m=\sum_l r_l e^{il\theta}$, and
$f(e^{i\theta})^m=e^{im\theta}(1+i+2\cos\theta)^m$; multiplying by $e^{-il\theta}$ and
integrating, $r_l=\frac1{2\pi}\int e^{i(m-l)\theta}(1+i+2\cos\theta)^m d\theta$, which is
\eqref{eq:phiint} with $j=m-l$. Finally $1+i+2\cos\theta=2(\cos\theta-c_0)$ with
$c_0=-(1+i)/2$, and $\cos j\theta=T_j(\cos\theta)$ identifies the Chebyshev coefficient.
\end{proof}

In these coordinates $N_l=|\phi_j|^2$ ($j=m-l$), the consecutive--ratio is
$\mu_l=N_l/N_{l-1}=|\phi_j|^2/|\phi_{j+1}|^2$, and $\zeta_l=r_{l-1}/r_l=\phi_{j+1}/\phi_j$.
Dividing \eqref{eq:rrec} and relabelling gives the clean \emph{outward} recurrence (verified
numerically for $m\le 14$ and algebraically equivalent to \eqref{eq:rrec}):
\begin{equation}\label{eq:star}
(m+j+1)\,\phi_{j+1}=(m-j+1)\,\phi_{j-1}-(1+i)\,j\,\phi_j .
\end{equation}

\section{The law is a clockwise turning, and is irreducible to magnitudes}
Set $\eta_j:=\phi_{j+1}/\phi_j$ and $X_j:=\Im(\overline{\phi_j}\phi_{j+1})=|\phi_j|^2\Im\eta_j$.
Since $\phi_{j+1}\overline{\phi_j}=r_{l-1}\overline{r_l}=P_l-iQ_l$ (with $l=m-j$), we have
$X_j=-Q_l$. Hence:
\begin{proposition}\label{prop:law}
The two--row law is exactly $X_j<0$ (equivalently $\Im\eta_j<0$) for all $1\le j\le m-1$:
the planar sequence $\phi_0,\phi_1,\dots,\phi_{m-1}$ turns strictly clockwise.
\end{proposition}

\begin{proposition}[Sign recurrence; irreducibility]\label{prop:sign}
Taking $\Im$ of \eqref{eq:star} multiplied by $\overline{\phi_j}$,
\begin{equation}\label{eq:Xrec}
(m+j+1)\,X_j+(m-j+1)\,X_{j-1}=-\,j\,|\phi_j|^2<0 \qquad(j\ge 1).
\end{equation}
Given $X_{j-1}<0$, \eqref{eq:Xrec} yields $X_j<0$ iff
$(m-j+1)|X_{j-1}|<j|\phi_j|^2$, which unwinds exactly to ``$Q_l>0$'', i.e.\ to the very
claim. Thus the sign data is self--referential: the law cannot be propagated from sign
information alone. Breaking the loop requires an external bound on the \emph{magnitude}
growth of $|\phi_j|$ --- precisely the role of the norm--ratio bound (``Lemma~1'').
\end{proposition}
This is, I believe, the cleanest explanation of the long--observed ``Cauchy--Schwarz is
circular'' obstruction. Verified ($m=14$): \eqref{eq:Xrec} holds and all $X_j<0$.

\section{Recessive solution and the continued fraction}
The recurrence \eqref{eq:star} (equivalently the general $(m+j+1)\phi_{j+1}=(m-j+1)\phi_{j-1}
+2c_0 j\phi_j$ for the Chebyshev coefficients of $(c-c_0)^m$) has a two--dimensional solution
space. Our $\phi_j$ is the \emph{recessive} (minimal) solution: the forward recurrence is
numerically unstable, amplifying the dominant companion solution. By Pincherle's theorem the
ratio of consecutive terms of the recessive solution is the value of a convergent continued
fraction, computed by the \emph{backward} (stable) recurrence.

\begin{theorem}[Continued--fraction form of the law]\label{thm:cf}
Let $\xi_j:=(m+j+1)\eta_j=(m+j+1)\phi_{j+1}/\phi_j$. Then
\begin{equation}\label{eq:cf}
\xi_{j-1}=\cfrac{(m+j)(m-j+1)}{(1+i)\,j+\xi_j},\qquad \xi_m=0,
\end{equation}
a terminating continued fraction, and $\eta_j=\xi_j/(m+j+1)$. The two--row law is exactly
\[
\Im\xi_j<0\quad\text{for all }1\le j\le m-1,
\]
i.e.\ every tail of the continued fraction \eqref{eq:cf} lies in the open lower half--plane.
\end{theorem}
\begin{proof}
From \eqref{eq:star}, dividing by $\phi_j$, $(m+j+1)\eta_j=(m-j+1)/\eta_{j-1}-(1+i)j$, so
with $\xi_j:=(m+j+1)\eta_j$ we get $(m-j+1)/\eta_{j-1}=\xi_j+(1+i)j$, hence
$\eta_{j-1}=(m-j+1)\big/\big[(1+i)j+\xi_j\big]$. Multiplying by $(m+j)$ yields
$\xi_{j-1}=(m+j)\eta_{j-1}=(m+j)(m-j+1)\big/\big[(1+i)j+\xi_j\big]$, which is \eqref{eq:cf}. The terminal value: $\phi_{m+1}=0$ gives $\eta_m=0$, so
$\xi_m=0$; the first backward step reproduces $\eta_{m-1}=\phi_m/\phi_{m-1}=(1-i)/(2m)$, the
known seed. The equivalence of the law to $\Im\xi_j<0$ is Proposition~\ref{prop:law} together
with $\Im\eta_j=\Im\xi_j/(m+j+1)$ and $m+j+1>0$.
\end{proof}
Verified exactly ($m=20$): \eqref{eq:cf} reproduces $\eta_j$ to machine precision and all
$\Im\xi_j<0$.

\begin{remark}[A general phenomenon]
For any $c_0$ with $\Im c_0<0$ and $\Re c_0\ne0$, numerical evidence (direct quadrature,
$m\le 15$, many $c_0$) indicates the Chebyshev coefficients of $(\cos\theta-c_0)^m$ turn
clockwise: $\Im(\overline{\phi_j}\phi_{j+1})<0$ for $j\ge1$. The hypothesis $\Re c_0\ne0$ is
necessary: at $c_0=-i$ the turning fails. Our case $c_0=-(1+i)/2$ sits on the locus
$|c_0|^2=\tfrac12$ (where the unrequired boundary value $X_0$ vanishes). A proof of this
general statement would prove the law for the whole family.
\end{remark}

\section{Why the value--region route, and why the box route fails}
The continued fraction \eqref{eq:cf} has the standard form $K(a_n/b_n)$ with
\[
a_n=(m+n)(m-n+1)>0\ (1\le n\le m),\qquad b_n=(1+i)\,n,\quad \arg b_n=\tfrac\pi4 .
\]
A \emph{value region} is a set $V$ with $s_n(V)\subseteq V$ for all $n$, where
$s_n(w)=a_n/(b_n+w)$; if $V\subseteq\{\Im w<0\}$ and the seeds enter $V$, all tails lie in
$\{\Im w<0\}$ and the law follows. A direct computation gives, for $a_n>0$,
\[
\Im s_n(w)=-\,a_n\,\frac{\Im b_n+\Im w}{|b_n+w|^2}=-\,a_n\,\frac{n+\Im w}{|b_n+w|^2},
\]
so $s_n$ maps $\{\Im w>-n\}$ into the open lower half--plane. The lower half--plane itself is
therefore \emph{not} invariant (it fails where $\Im w<-n$); the obstruction is concentrated at
the innermost, small--$n$ elements. The remaining task is exactly to produce an invariant
sub--region $V\subseteq\{\Im w<0\}$ (a disk or lens), which is the classical content of the
parabola / Thron--Jones value--region theorems for continued fractions with positive partial
numerators and partial denominators in a fixed half--plane.

\paragraph{The axis--aligned envelope (``box'') programme is inadequate.} The earlier
reduction bounded $\mu_\ell\le U_j=(2m+3j+3)/(2m-3j)$ by propagating a rectangle
$[L_j,U_j]\times[T_j,\bar T_j]$ in $(\mu_\ell,t_\ell)$ and checking four corner inequalities,
with $L_j,T_j,\bar T_j$ taken from the fixed--$j$ asymptotics
$\mu_\ell=1+\tfrac{3(2j+1)}{2m}+\tfrac{36j^2+16j-1}{8m^2}+\cdots$. I tested this directly:
it fails, and the failure is structural. The fixed--$j$ expansion has effective parameter
$j^2/m$ (successive terms scale as $(j^2/m)^k$), so it is valid only for $j\lesssim\sqrt m$;
but the band is $2\le j<2m/3$, which reaches $j\sim m$. On the bulk of the band the
``envelopes'' are simply wrong functions. (The hard end is small $j$, where the slack
$U_j-\mu_\ell=\tfrac{2j+1}{8m^2}+\cdots$ is $O(1/m^2)$ and the expansion \emph{is} valid;
the large--$j$ end has $O(1)$ slack and is easy.) Moreover the box decorrelates $(\mu,t)$,
which are tied along the orbit ($\mu-1\approx 6t$ to leading order); a correlated tube, not a
rectangle, is required. Both defects are absent in the continued--fraction formulation, where
the single complex variable $\xi_j$ carries the correlation automatically and the
value--region condition is uniform in $j$.

\section{The remaining gap, precisely}
\begin{conjecture}[Value region in the lower half--plane]
There is a region $V\subseteq\{\Im w<0\}$, e.g.\ a disk or half--plane depending on $(m,j)$,
with $s_n(\,\overline V\,)\subseteq V$ for the maps $s_n(w)=\dfrac{(m+n)(m-n+1)}{(1+i)n+w}$,
into which the terminal data $\xi_m=0$ feeds. Establishing this proves $\Im\xi_j<0$ for all
$1\le j\le m-1$, i.e.\ the entire two--row $d{=}4$ law, without recourse to the
$O(1/m^2)$--tight norm--ratio bound.
\end{conjecture}

\section{Verification}
All identities verified exactly in $\Z[i]$/\texttt{Fraction} (scripts in
\texttt{scratch/2026-06-05-lemma1/}): \eqref{eq:phiint} ($m\le12$), \eqref{eq:star} (all
$j$, $m\le14$), the sign recurrence \eqref{eq:Xrec} with $X_j<0$ ($m\le14$), and the
continued fraction \eqref{eq:cf} reproducing $\eta_j$ to $10^{-16}$ with all $\Im\xi_j<0$
($m=20$). The box--programme failure and the $j^2/m$ expansion scale were confirmed over
$6\le m\le 3000$.

\section*{Summary of what is proved here}
\begin{itemize}
\item[\textbf{R1}] (Lemma~\ref{lem:cheb}) $\phi_j=r_{m-j}$ are Chebyshev coefficients of the
pure power $(2c+1+i)^m$; clean integral form \eqref{eq:phiint} and outward recurrence
\eqref{eq:star}.
\item[\textbf{R2}] (Props.~\ref{prop:law},~\ref{prop:sign}) The law is the clockwise turning
$X_j<0$, and the sign recurrence \eqref{eq:Xrec} shows this is \emph{irreducible} to sign
data: a magnitude input is genuinely necessary.
\item[\textbf{R3}] (Theorem~\ref{thm:cf}) $\phi_j$ is the recessive solution; the law is
``all tails of the continued fraction \eqref{eq:cf} lie in the lower half--plane.''
\item[\textbf{R4}] The axis--aligned envelope programme is structurally inadequate (invalid
expansion for $j\gtrsim\sqrt m$; decorrelation); the value--region route is uniform and is
the natural classical setting.
\end{itemize}
\noindent\textbf{Remaining gap (single, classical):} a lower--half--plane value region for
the continued fraction \eqref{eq:cf}.
\end{document}
